\begin{abstract}
Selecting an action chunk from a visual policy requires predicting consequences that distinguish closely related proposals. Predicting spatial features retains a large visual target, whereas directly regressing a score provides no explicit future representation. We investigate Conditional Trajectory Abstraction (CTA), a learned prediction target for this decision: a source encoder compresses sampled future observations jointly with observed context; a world model forecasts that code from context and an action chunk; and a goal-conditioned reader scores the forecast without access to proposed actions. Ranking and feature-reconstruction supervision shape the source, followed by categorical prediction and reader-score consistency. Deployment uses 48 continuous expected-quantization coordinates, rather than future images. On identical eight-candidate PushT development banks, CTA recovers 0.680 of the available geometry improvement, compared with 0.635 for endpoint prediction and 0.502 for a parameter-matched direct scorer; paired 95\% intervals support both ranking differences. Historical closed-loop results are less decisive: CTA succeeds in 140/200 episodes versus 138/200 for endpoint prediction, and proposal batch sensitivity limits reproducibility. Longer-horizon diagnostics also expose a source-to-prediction gap. These findings support compact future-supervised targets for action ranking while identifying prediction fidelity, native-objective alignment, and reproducible held-out control as the outstanding tests of their practical value.
\end{abstract}
